%% 第 1 章  自动生成于 scripts/tex/md2tex.py，请勿直接手改
\chapter{绪论}
\label{chap:1}

\section{从内核内存安全缺陷到持续随机化}

内核的内存安全缺陷难以被完全消除，运行时防护因此不可或缺。内核运行于最高特权级，管理页表与物理内存，一旦被攻陷，攻击者即可取得系统级控制权。2024 年 1 月公开的 CVE-2024-1086 是一个典型例子：该释放后使用缺陷位于网络过滤子系统 \lstinline[breaklines=false]@nf_tables@，进入主线约十年才被发现，公开两个月后利用代码出现，随后被美国网络安全和基础设施安全局（CISA）列入已知被利用漏洞目录，并被确认用于勒索软件攻击。

这并非孤例：Linux 内核社区自 2024 年 2 月起成为 CVE 编号机构，截至 2026 年 9 月已为 10\,705 条缺陷分配编号。表~\ref{tab:1-1} 按缺陷类型统计了这批 CVE。本文把缺陷分为三类：时间类缺陷指针指向的对象已经失效，包括使用释放内存、空指针解引用、内存释放缺失与多重释放；空间类缺陷指针越出了对象的边界，包括越界读、越界写、整型溢出与数组越界；其余归入非内存类，包括条件竞争、锁使用不当与初始化错误。前两类合称内存安全缺陷，合计占全部 CVE 的 81.3\%，其中时间类 46.6\%、空间类 34.7\%。

这内存的时间错误与空间错误正是后续攻击步骤的主要入口：时间类对应图~\ref{fig:1-1} 中的指针悬空，空间类对应指针越界。缺陷持续出现的根源在于内核以不提供内存安全保证的 C 语言写成；Rust 正在进入内核，但只覆盖新增代码\citep{peng2025asterinas}，千万行存量 C 代码短期内无法替换。动态检测工具能在测试阶段发现其中一部分缺陷，但插桩开销使其难以在生产内核中长期开启\citep{cho2023bokasan}。由缺陷到利用原语的转化在内核上已相当成熟\citep{wu2018fuze,chen2020koobe,zeng2022playing,lee2023pspray,liang2024kleak}。因此，本文假定攻击者已经取得内存读写能力，考察如何阻止其进一步获取代码地址并劫持控制流。

\begin{table}[!htbp]
    \bicaption{\enspace Linux 内核 CVE 的缺陷类型分布}
    {\enspace Distribution of Defect Types in Linux Kernel CVEs (2024.1—2026.9)}
    \label{tab:1-1}
    \centering
    \scriptsize
    \setlength{\tabcolsep}{5pt}
    \renewcommand{\arraystretch}{1.05}
    \begin{tabular}{llrrrrr}
        \hline
        类别 & 缺陷类型 & 2024 & 2025 & 2026\textsuperscript{*} & 合计 & 占比/\% \\
        \hline
        \multirow{5}{*}{内存时间漏洞}
          & 使用释放内存 & 552 & 490 & 1\,065 & 2\,107 & 19.7 \\
          & 空指针解引用 & 357 & 349 & 347 & 1\,053 & 9.8 \\
          & 内存释放缺失 & 450 & 347 & 820 & 1\,617 & 15.1 \\
          & 多重释放 & 61 & 55 & 98 & 214 & 2.0 \\
        \cline{2-7}
          & 小计 & 1\,420 & 1\,241 & 2\,330 & 4\,991 & 46.6 \\
        \hline
        \multirow{5}{*}{内存空间漏洞}
          & 越界读 & 372 & 321 & 621 & 1\,314 & 12.3 \\
          & 越界写 & 318 & 286 & 553 & 1\,157 & 10.8 \\
          & 整型溢出 & 189 & 163 & 316 & 668 & 6.2 \\
          & 数组越界 & 168 & 145 & 260 & 573 & 5.4 \\
        \cline{2-7}
          & 小计 & 1\,047 & 915 & 1\,750 & 3\,712 & 34.7 \\
        \hline
        \multirow{4}{*}{非内存安全漏洞}
          & 条件竞争 & 268 & 231 & 402 & 901 & 8.4 \\
          & 锁使用不当 & 186 & 159 & 278 & 623 & 5.8 \\
          & 初始化错误 & 143 & 122 & 213 & 478 & 4.5 \\
        \cline{2-7}
          & 小计 & 597 & 512 & 893 & 2\,002 & 18.7 \\
        \hline
        \multicolumn{2}{l}{\textbf{内存安全漏洞合计}} & \textbf{2\,467} & \textbf{2\,156} & \textbf{4\,080} & \textbf{8\,703} & \textbf{81.3} \\
        \multicolumn{2}{l}{CVE 总数} & 3\,064 & 2\,668 & 4\,973 & 10\,705 & 100.0 \\
        \hline
    \end{tabular}
    \fignote{\textsuperscript{*}2026 年为部分年份，数据截至 2026 年 8 月 8 日的仓库快照，与表~\ref{tab:1-2} 同一份快照。}
\end{table}

内存错误本身并不直接构成危害，危害来自攻击者在其之上构造的利用链。图~\ref{fig:1-1} 给出内存安全的攻击模型：先使指针越界或悬空，再经该指针实施非法读写，进而篡改代码、代码指针、数据变量或输出数据，结果分为代码破坏、控制流劫持、纯数据攻击和信息泄露四类。防御技术沿这条链演进。NX 与 DEP 以页表权限区分可写页与可执行页，切断了由篡改代码到执行注入代码的直接路径；攻击者转而复用程序中已有的指令片段，以返回地址、函数指针或跳转表为落脚点把 gadget 串联成链，ROP\citep{shacham2007geometry}、JOP 与 COOP 分别以返回指令、间接跳转与虚函数调用组织片段。这类攻击的共同前提是确定代码的位置。地址空间布局随机化（ASLR）正作用于这一前提：Linux 3.14 引入的 KASLR 随机化内核代码段的基址，FG-KASLR\citep{kernel2021fgkaslr}进一步提出在启动时按函数重排。随机化无须针对每一种 gadget 组织方式分别定义规则，即可同时提高多类代码重用攻击的难度。

\begin{figure}[!htbp]
    \centering
    \includegraphics[width=0.88\textwidth,height=0.40\textheight,keepaspectratio]{schematic/ch1-memory-error-model}
    \bicaption{\enspace 内存错误利用模型及不同防御策略的作用阶段}
    {\enspace Memory-Error Exploitation Model and the Stages at Which Defenses Act}
    \label{fig:1-1}
\end{figure}

本文只防御由代码位置泄露通往控制流劫持的这条路径。图~\ref{fig:1-1} 中的着色部分即本文涉及的攻击路径：代码指针篡改、代码及代码索引读取，以及由此形成的控制流劫持。数据指针篡改与纯数据攻击不在直接防御范围之内：前者被视为攻击者已具备的读写原语；后者无须获知 gadget 地址，也不产生异常的控制流，例如直接改写进程凭据结构中的 \lstinline[breaklines=false]@uid@ 字段，需要数据流完整性或内存安全强制机制来防护，二者可与本文方法配合使用。第 2.6 节给出完整的威胁模型与防御边界。

内核部署防御机制受两条用户态没有的约束。其一，内核缺少同一软件栈内更高特权级的强制执行者：用户态的权限边界由内核配置页表并由处理器实施，而内核自己掌管页表，SMEP、SMAP、PXN、PAN 与代码段只读等机制的有效性都以其控制状态不被篡改为前提；要形成独立于内核的强制边界，只能借助不由普通页表直接控制的部件，例如以每处理器核寄存器保存权限状态的管理态保护键、只对数据访问触发的处理器调试观察点，或更高特权级的地址翻译。其二，内核能容忍的开销与兼容性代价更低：它位于系统调用、中断与异常的关键路径上，附加开销会传导至整个系统，且必须控制中断延迟；同时内核包含大量遗留代码、手写汇编与运行时代码修改，依赖全局静态分析或大范围改变代码语义的方案难以覆盖所有执行路径。

在这两条约束下，随机化是三类代表性防御路线中最适合内核的一种。控制流完整性先确定允许的转移关系，再在间接转移处验证目标，其效果取决于控制流图的精度\citep{abadi2005cfi}；C 语言的函数指针与运行时注册使精确分析困难，硬件辅助的 CET-IBT、BTI 与 kCFI\citep{shanbhogue2019cet,arm2023armv8,clang2022kcfi}只能验证落地点标记或类型签名这类可快速判定的属性，带相同标记或原型的目标仍构成较大的等价类\citep{gaidis2023fineibt,burow2017cfi}；指针认证虽能以上下文区分调用点，实际系统中却出现数千个调用点共享一个上下文的情况\citep{liljestrand2019pac}。内存隔离先要确定保护对象，再描述合法访问规则，这条研究线已从进程内隔离发展到内核分区化\citep{vahldiek2019erim,hedayati2019hodor,dautenhahn2015nested,huang2022ksplit,guo2025bulkhead,lefeuvre2025sok}，但内核中可被利用的对象随攻击方法变化，合法访问路径又随运行状态改变，为一般对象构建完整规则需要很强的语义分析能力。随机化主要依赖代码及其地址引用关系，这些信息大多可由编译器、链接器与重定位记录提供，对程序语义的依赖远小于前两者；装载时随机化不在每次转移或访问处检查，运行期开销也较低。

随机化依赖代码布局在运行期间保密，而信息泄露正是针对这一前提。攻击者可以利用越界读直接读出内存，可以根据访问是否引发崩溃探测地址有效性，也可以借助微体系结构侧信道；针对内核的去随机化侧信道已形成一条持续的研究线，从地址翻译缓存时序\citep{hund2013practical}、DrK\citep{jang2016drk}与分支预测器冲突\citep{evtyushkin2016jump}，到页表隔离部署之后仍然有效的预取侧信道\citep{liu2023entrybleed,lipp2022amd}、内核数据结构侧信道\citep{maar2025kernelsnitch}与瞬态执行攻击\citep{canella2019transient}，随机化一侧也有抵抗侧信道的尝试\citep{gens2017lazarus}。JIT-ROP\citep{snow2013just}是典型的信息泄露与控制流劫持的攻击：从一个泄露的代码指针出发扫描可达代码页，在运行时搜集 gadget 并构造攻击链，对应图~\ref{fig:1-1} 中先经“用指针读”获取布局、再经“篡改代码指针”完成劫持的两段路径。装载时随机化只在启动时改变一次布局，此后不再变化，攻击者可以反复探测并逐步恢复整个映像的布局，KASLR 无法使已经泄露的信息失效。

持续随机化把泄露地址的有效期压到攻击完成时间之下，其强度由随机化间隔与覆盖范围两个参数决定。Shuffler\citep{williams2016shuffler}在用户态每隔 50~ms 重排一次代码，而当时报告的 JIT-ROP 攻击完成时间为 2.3--378~s；Ahmed 等人\citep{ahmed2020methodologies}测量 gadget 的收敛速度后估计，要使攻击者难以在一个周期内集齐图灵完备的 gadget 集合，间隔上界约为 1.5--3.5~s。间隔越短，泄露信息可用的时间越短；覆盖范围越大，布局长期不变的代码越少；二者都以代码索引维护、执行状态协调与系统兼容的成本为代价。现有方案通过限制频率或范围控制这一成本，面向全部内核代码的持续随机化仍缺少可部署的实现。下一节先说明现有方案把范围限定在可加载内核模块所留下的盲区有多大，再由此提出本文要解决的问题。

\section{研究问题与面临的挑战}

将持续随机化用于内核时，主要困难来自运行时成本和执行正确性。内核执行环境包含系统调用、中断和异常等多种上下文，对延迟、稳定性和性能开销均有较严格的要求。持续随机化的成本主要集中在代码索引更新和执行流处理两个方面。代码位置发生变化后，系统需要识别并修正仍然指向旧地址的引用；对于规模较大且广泛使用代码指针的内核，这一过程可能涉及大量索引。与此同时，在旧代码被替换或解除映射之前，系统还需确认没有执行流继续使用该代码。内核中的中断、异常及其他非抢占上下文无法通过普通线程调度机制直接暂停，这进一步增加了执行状态协调的难度。

受上述成本限制，已有内核持续随机化研究通常只覆盖部分内核代码。Adelie\citep{nikolaev2022adelie}、Chaos\citep{chang2024chaos}、Remix\citep{chen2016remix}和 Dbox\citep{li2024dbox}等方案主要将驱动程序或可加载内核模块（Loadable Kernel Module，LKM）作为随机化对象。这些组件中包含较多第三方代码，其开发和测试过程可能不同于核心内核，因此常被选作优先保护对象。

这种保护范围主要反映了工程实现上的取舍，不意味着 built-in 内核代码不需要防护。以 Adelie 为例，其作者指出，对整个内核实施重新随机化需要处理核心内核中的遗留代码和底层实现，工程规模较大，因此该方案将持续随机化集中于被认为更容易受到攻击的组件\citep{nikolaev2022adelie}。这一设计降低了实现和运行成本，但没有解决 built-in 代码的持续随机化问题。

评估这一范围划分是否合理，需要考察两个相互独立的方面：其一是内核代码的可模块化属性，即由构建体系决定的结构性边界；其二是漏洞在不同代码区域中的实际分布，即上述边界是否与安全风险相符。

由于现有方案以可加载内核模块为保护对象，一段代码能否纳入其保护范围，首先取决于该代码能否被编译为模块。代码的可模块化属性由 Linux 内核构建系统在源码层面确定，并不随具体部署配置而改变。

在 Linux 的 Kbuild 体系中，目标文件的编译方式由所在目录 Makefile 中的赋值形式及相应的 Kconfig 配置符号共同决定。\lstinline[breaklines=false]@obj-y@ 表示目标文件无条件编入内核映像；\lstinline[breaklines=false]@obj-$(CONFIG_X)@ 则根据配置符号 \lstinline[breaklines=false]@CONFIG_X@ 的取值决定编译方式。若该符号在 Kconfig 中声明为 \lstinline[breaklines=false]@bool@ 类型，其取值仅包括 \lstinline[breaklines=false]@y@ 和 \lstinline[breaklines=false]@n@，相应代码只能被编入内核映像或不参与编译；只有声明为 \lstinline[breaklines=false]@tristate@ 类型的符号才允许取值 \lstinline[breaklines=false]@m@，相应代码因而具备编译为模块的可能。

由此可以得到一个不依赖具体内核配置的判定标准：若某个源文件在所有合法配置下均不能被编译为模块，则该文件始终处于现有部分随机化方案的保护范围之外。对 Linux 内核源码树的分析表明，多个核心子系统中的大量代码均符合这一条件：

\begin{itemize}
    \item \lstinline[breaklines=false]@kernel/@ 目录的 Makefile 中没有 \lstinline[breaklines=false]@obj-m@ 条目。调度器（\lstinline[breaklines=false]@kernel/sched/@）、锁与 RCU（\lstinline[breaklines=false]@kernel/locking/@、\lstinline[breaklines=false]@kernel/rcu/@）、中断处理（\lstinline[breaklines=false]@kernel/irq/@）以及时间子系统（\lstinline[breaklines=false]@kernel/time/@）的相关目标文件，或者直接列入 \lstinline[breaklines=false]@obj-y@，或者由 \lstinline[breaklines=false]@bool@ 型配置符号控制，均不能被编译为模块。
    
    \item \lstinline[breaklines=false]@mm/@ 目录的 Makefile 中同样没有 \lstinline[breaklines=false]@obj-m@ 条目。页面分配器（\lstinline[breaklines=false]@page_alloc.c@）、slab 内存分配器（\lstinline[breaklines=false]@slub.c@、\lstinline[breaklines=false]@slab_common.c@）、页缓存（\lstinline[breaklines=false]@filemap.c@）以及内存映射（\lstinline[breaklines=false]@mmap.c@、\lstinline[breaklines=false]@memory.c@）等组件均列入 \lstinline[breaklines=false]@obj-y@。
    
    \item \lstinline[breaklines=false]@fs/@ 目录中，VFS 核心层与具体文件系统采用不同的构建方式。VFS 核心文件，包括 \lstinline[breaklines=false]@open.c@、\lstinline[breaklines=false]@read_write.c@、\lstinline[breaklines=false]@namei.c@、\lstinline[breaklines=false]@file_table.c@、\lstinline[breaklines=false]@dcache.c@ 和 \lstinline[breaklines=false]@inode.c@ 等，直接列入顶层 \lstinline[breaklines=false]@obj-y@；\lstinline[breaklines=false]@ext4/@、\lstinline[breaklines=false]@xfs/@ 和 \lstinline[breaklines=false]@btrfs/@ 等具体文件系统则受 \lstinline[breaklines=false]@tristate@ 型配置符号控制，可以被编译为模块。
\end{itemize}

上述构建关系表明，负责进程调度、内存管理和 VFS 核心功能的代码主要以内建形式存在。仅以可加载内核模块为保护对象的随机化方案，在机制上无法覆盖这些代码。

另一项检验依据是漏洞的实际分布。Linux 内核社区于 2024 年 2 月成为 CVE 编号机构（CNA）。此后由内核社区发布的 CVE 记录包含具体的受影响源文件，为按文件粒度分析漏洞分布提供了数据条件。

本文采用以下方法进行统计：

\begin{enumerate}
    \item 从内核社区维护的漏洞信息仓库中提取编号年份为 2024、2025 和 2026 的已发布 CVE 记录，并解析每条记录所列的受影响源文件。
    
    \item 从受影响文件所在目录开始，沿 Kbuild 目录层级向上追溯其构建依赖，确定控制该目标文件及其上级目录的赋值语句和配置符号。直接列入 \lstinline[breaklines=false]@obj-y@，且构建依赖链仅包含 \lstinline[breaklines=false]@bool@ 型配置符号的文件，被判定为只能内建；若依赖链中存在 \lstinline[breaklines=false]@tristate@ 型配置符号，则该文件被判定为可模块化。
    
    \item 对于涉及多个源文件的 CVE，只要其中一个受影响文件被判定为只能内建，便将该 CVE 归入“只能内建”类别。该判定表示漏洞涉及的代码至少有一部分不能被现有部分随机化方案覆盖。
\end{enumerate}

上述判定依据源码中的 Kbuild 与 Kconfig 定义，不以特定发行版的内核配置为前提。2024 年 1 月至 2026 年 8 月的统计结果见表~\ref{tab:1-2}。

\begin{table}[!htbp]
    \bicaption{\enspace Linux 内核 CVE 所涉代码的可模块化性}{\enspace Modularizability of the Code Involved in Linux Kernel CVEs (2024.1—2026.8)}
    \label{tab:1-2}
    \centering
    \footnotesize
    \setlength{\tabcolsep}{3pt}
    \renewcommand{\arraystretch}{1.25}
    \begin{tabular}{lrrrr}
        \hline
        编号年份 & CVE 总数 & 只能内建 & 可模块化 & 未判定\\
        \hline
        2024 & 3 064 & 713（23.3\%） & 1 949（63.6\%） & 402（13.1\%）\\
        2025 & 2 668 & 610（22.9\%） & 1 778（66.6\%） & 280（10.5\%）\\
        2026* & 4 973 & 987（19.8\%） & 3 568（71.7\%） & 418（8.4\%）\\
        合计 & 10 705 & 2 310（21.6\%） & 7 295（68.1\%） & 1 100（10.3\%）\\
        \hline
    \end{tabular}
    \fignote{2026 年为部分年份，数据截至 2026 年 8 月 8 日的仓库快照。}
\end{table}

统计期内共收集到 10 705 个 Linux 内核 CVE，其中 2 310 个涉及在所有内核配置下均无法编译为模块的代码，占总数的 21.6\%。若排除构建方式无法确定的记录，在其余 9 605 个 CVE 中，该比例为 24.0\%。对于这些 CVE，Adelie、Chaos 和 Dbox 等仅保护模块或驱动的方案无法覆盖其中的内建代码。该缺口由代码的构建属性决定，不能通过调整内核配置消除。

将构建方式与缺陷类型交叉统计后可以看到，在 8 703 个内存安全缺陷 CVE 中，有 1 806 个涉及只能内建的代码，占 20.8\%。这两类缺陷可为越界访问、释放后使用及后续控制流劫持提供基础原语，而现有部分随机化方案无法保护其中以内建形式存在的代码。因此，将随机化范围限定于模块和驱动，会遗漏一部分与控制流劫持直接相关的漏洞。

按子系统看，只能内建的 CVE 主要分布在 \lstinline[breaklines=false]@kernel/@（413 条）、\lstinline[breaklines=false]@mm/@（207 条）、\lstinline[breaklines=false]@net/core/@（121 条）、\lstinline[breaklines=false]@net/ipv4/@（109 条）、\lstinline[breaklines=false]@arch/x86/@（84 条）、\lstinline[breaklines=false]@security/@（74 条）、\lstinline[breaklines=false]@io_uring/@（61 条）与 \lstinline[breaklines=false]@block/@（60 条），正是本节开头所列举的那些核心子系统。

上述分析表明，有必要将 built-in 内核代码纳入持续随机化的保护范围。这些代码承担着内核的基础功能，并广泛分布于系统的高频执行路径。例如，VFS 核心文件 \lstinline[breaklines=false]@fs/open.c@ 和 \lstinline[breaklines=false]@fs/read_write.c@ 包含大量文件相关系统调用的公共处理逻辑；调度器管理系统中的进程执行；内存分配器则由各内核组件共同使用。这些区域中的漏洞一旦被利用，可能影响系统的整体安全性与稳定性。因此，如何覆盖 built-in 代码构成了本文研究的基本问题。

持续随机化需要回答三个问题：如何随机化 built-in 代码，何时随机化其中的哪些函数，以及随机化区域之外的代码如何保护。研究内容一、二、三依次回答这三个问题。随机化机制是后两者的运行基础，触发策略决定随机化区域在运行期如何伸缩，边界保护则面对任何部分随机化都会留下的非随机化区域。三项内容各自成立：第 4 章与第 5 章所依赖的随机化原语在文献中已有现成方案，本文的原型由第 3 章提供。

当保护范围由内核模块扩展至 built-in 代码后，上述三个问题分别面临新的技术约束，同时受到一项共同问题的影响。该问题决定了三项研究内容之间的接口与配合方式，因而有必要首先加以分析。

\subsection{共性挑战：built-in 代码缺乏明确的随机化边界}

现有部分随机化方法利用驱动与核心内核之间的接口划分随机化区域。驱动只经导出符号表调用内核，内核只经注册的操作表回调驱动，全部跨界引用可以收拢进一张表，跨边界控制流集中于数量有限的入口，因而可以通过监控这些入口追踪执行状态并验证控制流转移。

Built-in 代码之间缺少可供利用的边界。核心子系统的调用关系相互交织，入口随执行上下文变化。调度器既可能从系统调用路径进入，也可能由中断触发；VFS 在处理文件操作时进入具体文件系统，并可能继续调用块层或网络层；内存分配器可由几乎所有内核组件调用，调用者可能处于进程、中断或原子上下文。按照源码目录或功能子系统划分随机化区域，得不到封闭且稳定的控制流边界。

边界不明确会影响持续随机化的三个环节。随机化机制需要知道哪些引用指向被搬移的代码，以及哪些执行流仍停留在其中；触发策略需要知道随机化区域的当前地址存放在何处，才能观测对它的读取；边界保护需要一个所有跨区域控制流都经过的位置，才能施加验证。

本文以跳板函数统一处理这一问题。每个被随机化的函数对应一个位置固定的入口跳板，被随机化代码调用的每个外部函数对应一个出口跳板，所有跨区域控制流都经过跳板；函数体的当前地址只存在于一条直接分支指令的立即数中，本文称这条指令为跳转桩，随机化时唯一改写的就是它。由此形成的边界具有三种用途：集中维护与代码位置有关的索引，观测进入和离开随机化区域的执行流，以及施加跨区域控制流验证。第 3 章给出跳板机制及其随机化实现，第 4 章在跳转桩区上施加只执行保护以识别泄露，第 5 章在跳板处施加分区验证。三项研究内容由同一边界机制衔接。

\subsection{挑战一：代码索引更新与执行流协调}

代码随机化会改变函数及指令的位置，所有依赖原代码地址的信息都必须随之调整，包括代码段中的直接调用与跳转、数据段中的函数地址、栈中的返回地址，以及全局偏移表等结构中的相关条目。本文将这类与代码位置相关、需要在随机化过程中维护的信息统称为代码索引（code index）。在 built-in 代码上处理代码索引，受到更新规模、地址语义、返回地址与执行流协调四方面的限制。

传统重定位方法需要修正代码段中的直接调用与跳转指令，追踪数据段中的函数地址，对活动线程进行栈展开并重定向返回地址，同时重新解析符号及更新相关链接表。若以需要处理的代码引用数量为 \(n\)，一次随机化的更新开销为 \(O(n)\)。built-in 代码之间的引用密集，且散布在代码、数据与 per-CPU 区域中，全量更新难以承受高频随机化。

地址语义进一步增加了索引维护的难度。内核 C 代码允许在指针与整数之间转换，函数地址也可能参与比较、哈希或数值运算。以保存函数地址的变量 \lstinline[breaklines=false]@x@ 为例，当该变量被当作数值使用时，随机化不应改变其数值语义；当它被解引用或用于控制流转移时，其中的地址又必须指向函数的新位置。因此，仅凭赋值操作无法确定该值后续将以何种语义使用。在赋值阶段将函数地址替换为跳转表下标会改变原有数值语义；若推迟到使用阶段再进行地址转换，则需要在运行时查询映射关系，并覆盖内核中多种形式的间接分支目标，映射规模与查询开销都会随之增加。

栈上的返回地址作为代码索引同样需要处理。它由处理器在调用时自动写入，指向函数体内部的当前位置。一个随机化函数在调用其它函数期间，其返回地址就在栈上，内核栈因而成为一份随时可读的布局副本；若该函数在此期间被搬移，返回时便落到已失效的位置。逐条内核栈展开并改写返回地址，又与更新规模的约束相冲突。

执行流协调是代码更新面临的另一项约束。随机化期间必须保证没有执行流仍在待移动的代码中，但核心内核路径包含大量非抢占上下文，如中断处理程序、持有自旋锁的临界区、RCU 读端临界区和显式关闭抢占的代码段。基于调度器的暂停方法不能直接停止这些执行流。例如，Adelie 使用 \lstinline[breaklines=false]@schedule_on_each_cpu()@ 并等待 RCU 宽限期，但 RCU 宽限期主要用于等待既有读端临界区结束，不能单独确认中断处理程序或持锁代码已经离开随机化区域。强制暂停持有自旋锁的执行流还可能使其他 CPU 长时间等待该锁，甚至形成死锁。 已有的用户态持续随机化方案以保留旧代码副本回避这一约束，在内核中旧副本意味着已泄露的地址在切换之后仍然可用。

\subsection{挑战二：随机化的触发时机与作用对象}

现有方案多采用固定时间间隔触发随机化。缩短间隔可以压缩地址泄露后的利用窗口，但会增加随机化次数和运行时开销；延长间隔虽然降低了开销，却为攻击者保留了更长的攻击时间。第 1.1 节引用的实验结果将有效随机化间隔限定在秒级，而内核全量随机化的单次成本使这一频率难以长期维持。

随机化的范围同样受制于开销。以跳板为间接层的方案中，一次调用只要有一端落在随机化区域内就要经过跳板：随机化函数占全部函数的比例为 \(p\) 时，经过跳板的调用占全部调用的比例为 \(2p-p^{2}\)，小范围时约为 \(p\) 的两倍。第 3 章的实测给出了这一杠杆在真实内核上的形式：随机化不到 1\% 的函数时常驻开销为 1.98\%，扩到 40\% 的 built-in 代码时为 28.83\%。随机化因而只能限定在一个子集内，而信息泄露却可能发生在内核的任何位置。固定周期还会对未受探测的代码重复执行随机化。一次随机化只有发生在代码地址被探测之后、泄露地址被用于控制流劫持之前，才能直接破坏当前攻击所依赖的布局；对未被探测的代码以同等频率随机化，只产生开销而没有相称的安全收益。

触发式随机化需要一个直接来自信息泄露的信号。只执行内存能够把对代码的读取变为硬件事件，但现有的内核只执行内存方案都是防护式的。内核持续地读取自己的代码：动态跟踪设施在插入探测点时读取原指令，指令替换与静态键在启动期读取并改写代码，模块加载器与实时补丁在应用重定位时读取代码，性能分析工具在栈回溯与符号解析时读取指令流。防护式方案必须为这些合法读取者维护一份完备的豁免清单，漏掉任何一条都会使内核在该路径首次执行时崩溃。由此产生三项挑战：如何使只执行保护不依赖一份必须完备的豁免清单；如何把读取事件转化为可靠的触发判据，并使从触发到完成的延迟短于攻击者可利用该次读取的时间；如何使常驻开销与随机化范围的规模解耦，使候选范围可以扩大到全内核而开销不随之上升。

\subsection{挑战三：非随机化区域中指针认证上下文的传播}

受开销所限，现有的内核持续随机化方法全部只能做部分随机化。Adelie、Chaos 与 Dbox 以模块与驱动为界，触发式随机化以泄露事件为界，随机化区域之外的代码布局在运行期不变。攻击者一旦持有指向这些代码的合法指针值，它就长期可用，仍可利用其中的 gadget 构造 ROP 或 JOP 链。非随机化区域的保护因此是所有部分随机化方法共同面对的问题。

如第 1.1 节所述，硬件辅助 CFI 可以为非随机化区域提供补充保护。商用硬件机制中，ARM 指针认证是唯一能以软件选取的上下文区分调用点的一种，原理上可以为每个调用点取不同的值，但受上下文传播约束：多个调用点调用同一函数时必须共享上下文，该约束沿调用图传播，使原本独立的调用点逐渐归入同一上下文。在 iOS 13 中，约有 6\,500 个调用点共享同一上下文\citep{liljestrand2019pac}。共享范围扩大后，通过同一认证的合法目标随之增加，间接调用检查的区分能力下降。

上下文传播由高扇入（high fan-in）函数引起。此类函数可由多个调用点到达，把原本相互独立的认证上下文连成一片。内核 C 代码广泛使用指针别名、间接调用、函数指针数组和动态回调注册，静态分析难以完整恢复这些函数的调用来源。由此产生三项挑战：如何在不做全局指针分析的前提下识别高扇入函数并让其余调用点得到唯一上下文；如何避免在每一次间接调用上都施加验证，使开销随随机化范围可控；如何使保护机制不绑定某一种随机化实现，从而可以施加到任何一种部分随机化方法留下的非随机化区域上。

\section{研究内容与主要贡献}

针对上述挑战，本文提出并实现了一套面向 Linux 内核 built-in 代码的持续随机化系统。该系统以随机化区域与非随机化区域的划分为基础，跨区域的控制流被收束到跳板这一唯一出入口，本文的三项工作都建立在这一划分之上。系统的总体设计目标有三：

\begin{enumerate}
    \item 不引入新的硬件设计。如第 1.1 节所述，已有的硬件随机化方案都需要修改处理器，无法用于现役系统。本文只使用现役商用处理器已经提供的部件：x86 的管理态保护键、ARM 的硬件调试观察点与 ARM 指针认证。三者都被用在其原本用途之外的场合。保护键原本用于内核内部的数据隔离，本文用它把跳转桩区置为只执行；观察点原本用于调试，本文用它在无虚拟化层的条件下捕获对跳转桩区的读取；指针认证原本用于指针完整性，本文用它保护随机化区域之外的控制流。
    \item 不修改被保护子系统的源码语义。改动集中在编译器与随机化运行时，被随机化的内核子系统不需要为适配本文机制而改写代码逻辑。
    \item 开销可承受。在真实负载下的端到端开销控制在个位数百分比。
\end{enumerate}

本文的三项研究内容依次对应第 3、4、5 三章。三者由同一边界机制衔接，但各自可以独立成立：第 4 章与第 5 章各自只对随机化原语提出一项要求，该原语在文献中已有现成方案。图~\ref{fig:1-2} 以内核持续随机化为研究方向，自上而下列出由它分解出的三个问题、每个问题面临的关键挑战，以及与之对应的三项研究内容；虚线箭头表示后一项研究内容建立在前一项之上。

\begin{figure}[!htbp]
    \centering
    \includegraphics[width=0.95\textwidth,height=0.46\textheight,keepaspectratio]{schematic/ch1-research-outline}
    \bicaption{\enspace 本文的研究思路与三章分工}{\enspace The Research Rationale of This Dissertation and the Division of Labour across Chapters 3 to 5}
    \label{fig:1-2}
\end{figure}

\subsection{研究内容一：基于跳板函数的低开销持续随机化}

要回答的问题：如何随机化 built-in 代码。对应第 3 章。

跳板本身不是本文提出的。Oxymoron\citep{backes2014oxymoron}以跳板替换函数指针，Shuffler 以索引表实现间接，MARDU\citep{jelesnianski2020mardu}与 CodeArmor\citep{chen2017codearmor}同样引入了地址间接层，把跳板用作随机化时只需更新一处的手段是这一领域的通行做法。本文的贡献在于跳板集中的是什么以及跳板还能做什么。

本文把代码索引分层处理。函数粒度内以位置无关代码消除索引，函数体可以被搬到任意地址而无需修改其内部的任何一条指令；函数粒度间以固定位置的跳板集中索引，函数体的当前地址只存在于跳转桩这一条直接分支指令的立即数中，一次随机化的改写量等于函数个数而与调用点数量无关。函数指针的值恒为跳板地址，内核中把函数地址当作数值比较、哈希与反查的代码模式不受影响，数值还是地址的判定被推迟到了控制流转移发生的那一刻。栈上的返回地址由与随机化无关的返回标记取代，标记只含跳板序号与函数内偏移，经重入桩解析到当前变体，内核栈不再承载当前布局，陈旧返回不再发生。

本文还扩展了跳板的职责。已有工作把跳板当作纯粹的地址间接层，本文把它同时用作随机化区域边界上的执行流观测点：入口与出口计数精确给出当前正在区域内执行的执行流集合，据此确定安全的暂停点，随机化不必保留任何旧代码副本，切换完成的那一刻旧地址就真正过期；对中断、异常与关抢占区间这类无法被调度器暂停的上下文，采用推迟而非阻塞的策略。

主要贡献：

\begin{enumerate}
    \item 提出基于函数粒度位置无关代码与跳转桩的分层代码索引组织方法，使一次随机化的改写量与被引用次数无关，并以返回标记消除栈上的布局副本与陈旧返回；
    \item 提出函数指针语义分离方法，在不改变指针值的前提下保证所有跨区域控制流转移必经跳板；
    \item 提出基于跳板进出计数的精确执行流追踪与推迟随机化策略，实现无代码副本的随机化并覆盖非抢占上下文。在 483 个函数的随机化范围下，单次随机化的停顿窗口为 1.22~μs，等待执行流离开当前函数体的时间在 0.04 至 0.18~μs 之间；在与 Adelie、Chaos 和 Dbox 的同条件对比中，本文机制在各档随机化间隔下的开销均最低。
\end{enumerate}

\subsection{研究内容二：基于代码索引泄露的触发式随机化}

要回答的问题：何时随机化哪些函数。对应第 4 章。

本文的关键观察是：随机化区域的当前地址只存在于跳转桩这一处，对它的数据读取就是泄露本身，不需要任何代理。据此本文把只执行内存的保护对象由整个代码区收窄为跳转桩区，并把它从阻断机制改用为检测机制。正常执行到达函数体的路径全部由直接分支构成，只被取指；内核的合法读取者读的是函数体的指令而非跳转桩；落在跳转桩区上的任何数据读取因而都没有合法解释，机制不需要为任何读取者维护豁免清单。读取事件驱动被读函数进入随机化，随机化范围由编译期静态选定改为运行时由泄露事件触发。

把只执行内存用作检测而非阻断并非本文首创。Heisenbyte\citep{tang2015heisenbyte}与 NEAR\citep{werner2016near}采用允许读取但使读到的代码不可用的思路，SafeHidden\citep{wang2019safehidden}以探测检测驱动用户态数据区域的重随机化。本文的不同在于保护对象与响应方式：保护对象是跳转桩区而非代码本身，合法读取者与被监控对象由此天然分离；响应是一次针对性的随机化，不改变被读代码的任何语义；机制实现在内核态。

常驻开销与候选范围的解耦由两级跳板完成。候选函数的一级跳板在未激活态只是一条到跳转桩的直接分支，不做任何记账，一次往返约 5 个周期；被读取的函数才激活为完整的跳板，约 54 个周期。激活是一次 4 字节的原子写入。候选集合因而可以扩到全内核，每次调用只为被读过的函数付完整的代价。

主要贡献：

\begin{enumerate}
    \item 提出以跳转桩区的读取事件作为持续随机化的触发信号，使随机化由泄露驱动而非由时钟驱动，且不依赖任何豁免清单；
    \item 提出两级跳板与跳板激活机制，使常驻开销与随机化的候选范围解耦；
    \item 给出内核态只执行内存的两条实现路径，x86 平台基于管理态保护键，ARM 平台基于硬件调试观察点，两者都只作用于数据访问而对取指透明，放行一次读取只影响本核。单次检测开销在 ARM 观察点路径上为 1.32~μs、在 x86 保护键路径上为 0.96~μs，基础负载下内核自身的读取每秒触发 0.02 至 0.44 次，从检测到随机化完成的延迟在基础负载下中位为 12 至 21~μs、99 分位不超过 110~μs。
\end{enumerate}

\subsection{研究内容三：基于随机化辅助的 ARM 指针认证机制}

要回答的问题：非随机化区域怎么办。对应第 5 章。

本文的关键观察是：随机化区域内的目标地址不可预测，进入该区域的间接调用不需要指针认证，随机化因此恰好提供了打断上下文传播链所需的东西。把造成传播的高扇入函数移入随机化区域，链两侧的调用点即可各自建立基于调用点地址的唯一上下文。高扇入函数的识别在运行时完成：指针在识别阶段以类型上下文签名，归属由独立记录跟踪，一个函数第二次被不同的调用点调用即判为高扇入并移入随机化区域，单一归属的函数则在上下文绑定时重签为调用点唯一上下文，全程不需要全局指针分析。验证在跳板处分区施加：进入随机化区域只验范围，离开时强制认证与白名单检查。该机制对随机化原语只有一项要求，即被识别为高扇入的函数持续处于随机化状态，因而可以施加到任何一种部分随机化方法留下的非随机化区域上。

主要贡献：

\begin{enumerate}
    \item 提出以随机化打断指针认证上下文传播链的协同机制，使随机化与控制流完整性互为条件；
    \item 提出基于运行时归属跟踪的高扇入函数识别与上下文绑定方法，不依赖全局指针分析；
    \item 提出跳板处的分区验证，并说明该机制与具体随机化方法解耦。最大等价类由类型上下文下的 141 降到 57，91\% 的间接调用点合法目标数不超过 5，代价是把 233 个高扇入函数移入随机化区域；施加到 Chaos 之上，开销增量为 0.7\%。
\end{enumerate}

\subsection{三项工作的关系}

三项工作对应持续随机化在内核场景下的三个问题：研究内容一压低单次随机化的代价并覆盖 built-in 代码，使高频随机化成为可能；研究内容二取消固定频率，使随机化只在泄露发生时、只对被读过的函数发生；研究内容三承认部分随机化必然留下非随机化区域，以指针认证补齐随机化够不到的部分。三者通过跳板耦合：研究内容一构造它并把代码索引集中到跳转桩，研究内容二保护跳转桩并在跳板上切换激活状态，研究内容三在跳板处施加跨区域验证。每一章的问题都由上一章的边界逼出：跳板的常驻开销限定了随机化范围，触发式随机化因此只随机化被读过的函数，而随机化区域之外的部分交给控制流完整性。

本文在 Linux 6.8 内核与 LLVM 14 上实现了完整原型，主平台为 ARM Neoverse N2，副平台为 Intel Xeon Gold 6448Y。第 3、4、5 章各自给出机制级评估与安全性分析，与 Adelie、Chaos 和 Dbox 的同条件对比见第 3.5.4 节，第 5 章机制在 Chaos 上的移植验证见第 5.5.4 节。

\section{本文的安全性论证方式}

防御机制的有效性不像性能那样可以直接测量，如何论证是一个方法学问题。本文在此明确声明所采用的证据层级，以及每一层的适用范围——\textbf{不同性质需要不同强度的证据，把只能被测量的性质写成定理，反而会削弱论证。}

本文的安全性论证由四层证据构成：

\textbf{第一层，可机械验证的结构性质。} 少数性质可以被完全验证，本文对这类性质给出严格论证并附验证程序。核心是一条：编译产物中不存在绕过跳板函数的跨区域控制流转移。该性质可由编译期的静态检查与对内核映像的反汇编扫描共同确认，是本文全部安全论证的基础——后续各层的论证都以它成立为前提。该检查由编译器插件在编译期执行，并以对内核映像的反汇编扫描作为构建后的复核，实现见第 3.4.1 节。

\textbf{第二层，概率与熵的定量分析。} 随机化类防御的强度本质上是概率性的，因此以概率模型而非定理来刻画。本文为攻击者在给定探测次数下的成功概率建立模型，并遵循一条原则：模型中所有简化假设都必须朝有利于攻击者的方向取值，从而使所得概率成为安全性的下界而非估计值。每条假设在提出时都需说明它为何是保守的。认证码伪造这一项的概率界见第 5.6 节。这一做法与该领域的通行实践一致——例如 Adelie 通过地址空间熵的直接计算给出攻击成功概率的上界，并与 Shuffler、CodeArmor 作横向对比。

\textbf{第三层，攻击复现实验。} 检测覆盖率、误报漏报率、随机化触发的及时性，这些性质无法演绎，只能测量。本文以真实 Linux 内核提权漏洞为攻击\textbf{案例}，把其已公开的利用路径按“获取原语、泄露布局、收集串联、劫持控制流”四步分解，在启用与未启用各档机制下逐步骤观测拦截结果，记录攻击在哪一步被拦截、拦截前攻击者来得及读取多少代码片段、以及触发拦截所需的探测次数分布。需要说明攻击\textbf{案例}与攻击\textbf{驱动}的分工：案例是真实的，即真实漏洞与其真实的利用步骤；而落在防御范围内的步骤由内核内受控自测忠实复现该步骤的行为，获取原语这一步则按威胁模型直接赋予，本文不实现、也不运行任何向用户态暴露内核任意读写的武器化载体。这样做既守住不引入两用攻击代码的底线，又使被拦截行为的起点在每次实验中完全一致，从而可把机制的响应与利用本身的偶然性分开。第 4.5.4 节以内核内受控自测复现“泄露地址至解引用”这一步骤，给出各时延档位下的攻击成功率；第 5 章机制不涉及该步骤的复现，其有效性由等价类规模的度量给出。

\textbf{第四层，残余风险的显式声明。} 本文明确声明防御边界之外的部分，包括纯数据攻击、硬件侧信道、以及依赖硬件规范正确实现的假设。特别地，ARM 指针认证已被证明可通过推测执行侧信道在不触发崩溃的情况下被爆破\citep{ravichandran2022pacman}，因此研究内容三的安全性论证以“攻击者无法获得 PA 验证预言机”为前提，该前提及其脆弱性在第 5.6 节单独说明。

本文不采用形式化定理的呈现形式。原因是：本文的安全性质中，只有第一层可以被真正证明，第二层是概率陈述，第三层是经验测量；将后两者以定理形式呈现，会把“在给定假设下的度量结果”误示为“演绎结论”，一旦某条假设在答辩或后续研究中被质疑，整条论证链的可信度会一并受损。

\section{论文组织结构}

本文共分 6 章。每一章对应一项独立的工作，各章自足：各自给出研究动机与设计需求、系统设计、实现细节、机制级实验与安全性分析。跨越三项工作的内容，即统一的威胁模型与安全不变式，放在第 2 章。

\textbf{第 1 章 绪论}：从内存安全攻防的坐标系出发，说明随机化路线的位置与硬件 CFI、内存隔离、硬件随机化各自的边界；通过内核构建体系的结构性事实与 CVE 分布的统计事实，提出 built-in 内核代码的保护盲区问题；给出持续随机化在内核场景下的核心矛盾，分析一个共性挑战与三个专属挑战；概述研究内容、主要贡献与安全性论证方式。

\textbf{第 2 章 相关研究综述与威胁模型}：综述内存安全攻防方法的演进、内核随机化的三个发展阶段、持续随机化三个核心问题上的已有解法、XOM 技术与硬件 CFI 技术的现状，并指出各自的不足与本文的出发点；在此基础上给出本文统一的威胁模型，定义攻击者能力、四类攻击向量、信任基础与明确的排除项，并由威胁模型导出三条贯穿全文的安全不变式。\textbf{后续各章的安全性分析均引用本章定义的威胁模型与不变式。}

\textbf{第 3 章 基于跳板函数的低开销持续随机化方法}：研究内容一。本章首先分析代码索引的规模、返回地址的语义与执行流暂停三个层次上的困难，然后给出跳板函数的指令序列、函数粒度的位置无关代码、以位置无关返回标记替代裸返回地址，以及以进出计数刻画活跃执行流的安全暂停，随后给出 LLVM 编译器扩展与内核运行时的实现，最后评估随机化延迟、随机化核心子系统的开销、与已有工作的同条件对比，以及随机化范围与跳板调用占比的关系。\textbf{跳板函数作为后两章的共同基础设施，其完整设计在本章给出。}

\textbf{第 4 章 基于代码索引泄露的触发式随机化}：研究内容二。本章首先分析固定间隔触发策略的两难与保护整个代码区域的代价，然后给出把只执行内存由阻断机制改用为检测机制的核心思想、只保护跳转桩区这一收窄后的监控对象，以及未激活与激活两态的两级跳板，随后给出 x86 平台基于管理态保护键与 ARM 平台基于硬件调试观察点的两条实现路径，最后评估两条路径的单次检测开销、各类负载下的触发频率与延迟、泄露地址的有效期与利用竞态，以及两态下的单次调用开销。

\textbf{第 5 章 基于随机化辅助的 ARM 指针认证机制}：研究内容三。本章首先分析部分随机化留下的非随机化区域与指针认证的上下文传播问题，然后给出以类型上下文配合归属跟踪在运行时识别高扇入函数、把高扇入函数移入随机化区域以打断传播链、并在跳板处施加分区验证的设计，随后给出编译器插桩与运行时实现，最后以每个间接调用点的合法目标数评估精度，评估识别的收敛性，并把该机制移植到已有的部分随机化方法上。

\textbf{第 6 章 总结与展望}：总结全文工作与创新点，讨论局限性与残余风险，展望后续研究方向。

第 3、4、5 三章采用统一的内部结构：引言与本章导航 → 现有工作与研究动机（以设计需求清单收尾）→ 系统设计 → 系统实现 → 机制级实验评估 → 安全性分析 → 本章小结。各章的实验只回答“本章机制是否达到了本章开头列出的设计需求”，同一组基准测试数据不在多章之间重复报告。

